\newpage
\section{Nov 18 Notes}

\begin{remark}
    The principal ideal $(k)$ is exactly the set of all elements of $R$ divisible by $K$. Thus, if $R = F[x]$, where $F$ is a field, then for $k := x - a$, for $a \in F$, the principal ideal $(x-a)$ is exactly the set of polynomials where $x = a$ is a root; equivalently, the set of polynomials which are divisible by $(x - a)$.
    
    Note that $(0) = \{0\}$ and we call this the \textit{zero ideal}.
\end{remark}

\begin{defbox}
    A \bold{principal ideal domain} (abbreviated (\bold{PID})) is a ring which is an integral domain and every ideal in it is principal, i.e generated by a single element. 
\end{defbox}

\begin{example}
    $\ZZ$ is a PID. This is because ideals are in particular additive subgroups and the only subgroups of $\ZZ$ are $n\ZZ$ by homework, where $n \in \ZZ$. Note that $n\ZZ = (n)$.
\end{example}

\begin{thmbox}
    Let $F$ be a field. Then $F[x]$ is a PID. In fact, every nonzero ideal in $F[x]$ is generated by a unique monic polynomial.
\end{thmbox}

\begin{proof}
    Let $I$ be a nonzero ideal of $F[x]$. Let $d \in I$ be a nonzero polynomial of least degree; since $F$ is a field, we can make it monic by multiplying it with the inverse of the leading term. Let us thus assume that $d$ is monic. Since $d \in I$, $(d) \subseteq I$ since $I$ contains all multiples of its elements. On the other hand, take any $f \in I$ and divide by $d$ with remainder:
    \[ f = qd + r \]
    where $\deg(r) < \deg(d)$, i.e. $r = 0$ since $d$ is of the least degree. Then $r = f - qd \in I$, so by the minimality of $\deg(d)$ among all nonzero elements of $I$, $f \in (d)$, so $I = (d)$.

    To show uniqueness, if $e$ is a monic polynomial where $(e) = I = (d)$, then $d \mid e$ and $e \mid d$, so $d = pe$ and $e = qd$, hence $\deg(e) \leq \deg(d)$ and $\deg(d) \leq \deg(e)$, so $\deg(d) = \deg(e)$, and since $q$, $p$ are constants, and $d$ and $e$ are monic, $p = q = 1$, so $d = e$.
\end{proof}

We can also define ideals generated by subsets, just like with subgroups.

\begin{defbox}
    Let $R$ be a ring, $S \subseteq R$. The \bold{ideal generated by $S$} is the smallest (with respect to inclusion) ideal of $R$ containing $S$, denoted $(S)$. Just like with subgroups, ideals are closed under intersections, so $(S)$ is the intersection of all ideals of $R$ containing $S$.
\end{defbox}

When $S = \{k\}$ or $S = \{k_1, \dots, k_n\}$, we write $(k_1, \dots, k_n)$. Note that the principal ideal generated by $k$ is exactly $(\{k\})$.

\begin{example}
    For $a$, $b \in R$, where $R$ is commutative, $(a,b) = Ra + Rb := \{ua + rb : u,v \in R\}$. Indeed, note that $Ra + Rb$ is indeed an ideal (not necessarily the smallest), so $Ra + Rb \supseteq (a,b)$. The converse clearly holds because $(a,b)$ is closed under multiplication by $R$, so $Ra$, $Rb \subseteq (a,b)$, and $(a,b)$ is also closed under addition, leading to the result.
\end{example}

\begin{defbox}
    Let $F$ be a field and $f$, $g \in F[x]$, with at least one of $f$ or $g$ nonzero. The \bold{greatest common divisor (GCD)} of $f$ and $g$, denoted $\gcd(f,g)$, is the unique monic polynomial of highest degree which divides both $f$ and $g$. 
\end{defbox}

\begin{corrbox}
    Let $F$ be a field, $f$, and $g \in F[x]$ with at least one being nonzero. Then $\gcd(f,g)$ exists and is unique. In fact $(f,g) = (\gcd(f,g))$ and Bézout's Identity holds.
\end{corrbox}

\begin{proof}
    Let $d$ be the unique monic polynomial in $F[x]$ with $(f,g) = (d)$. Then $d \mid f$ and $d \mid g$, because $f$, $g \in (d)$. Also, $d = uf + vg \in (f,g)$, hence every common divisor of $f$ and $g$ also divides $d$, so $d$ has the highest common divisor of $f$ and $g$, thus $d = \gcd(f,g)$. 

    To show uniqueness, if $e$ is also a monic polynomial of highest degree which divides both $f$ and $g$, then we have that $(f,g) \subseteq (e)$, but also $e \mid d$, because every common divisor of $f$ and $g$ divides $d$, so $(e) \subseteq (d)$, hence $(e) = (d)$ by uniqueness of the previous theorem. 
\end{proof}

\begin{defbox}
    We call polynomials $f$ and $g$ in $F[x]$ \bold{coprime} if $\gcd(f,g) = 1$. 
\end{defbox}

\begin{thmbox}[Euclidean Algorithm]
    Let $F$ be a field and $f$, $g \in F[x]$, with $g \neq 0$. Set $r_0 := f$ and $r_1 := g$. Divide with remainder successively:
    \begin{align*}
        &r_0 = q_0r_1 + r_2 \hspace{0.5cm} \deg(r_2) < \deg(r_1) \implies (r_0, r_1) = (r_1, r_2) \\
        &r_1 = q_1r_2 + r_3 \hspace{0.5cm} \deg(r_3) < \deg(r_2) \implies (r_1, r_2) = (r_2, r_3) \\
        &\vdots \\
        &r_{k-1} = q_{k-1}r_{k} + r_{k+1} \hspace{0.5cm} (r_0, r_1) = (r_k)
    \end{align*}
    where $r_{k+1} = 0$, hence $\gcd(f,g) = r_k$.
\end{thmbox}

\begin{example}
    In $\QQ[x]$, let $f := x^5 - 1$, $g = x^3 - 1$. By the Euclidean Algorithm:
    \begin{align*}
        x^5 - 1 &= x^2(x^3 - 1) + (x^2 - 1) \\
        x^3 - 1 &= x(x^2 - 1) + (x - 1) \\
        x^2 - 1 &= (x+1)(x-1) + 0 
    \end{align*}
    Thus $\gcd(x^5 - 1, x^3 - 1) = x - 1$ and we can find the $u$ and $v$ satisfying Bézout's identity using back substitution:
    \[ x - 1 = (-x)(x^5 - 1) + (1 - x^2)(x^3 - 1) \]
\end{example}